% ЗАКЛЮЧЕНИЕ, без номер.
% Обобщава характеристиките на решението, предимствата и недостатъците му,
% резултатите и посоките за по-нататъшна работа.
\section*{ЗАКЛЮЧЕНИЕ}
\label{sec:conclusion}
\addcontentsline{toc}{section}{ЗАКЛЮЧЕНИЕ}

Проектът разработва система за подпомагане на решенията при избор и подреждане
на портфейл от ИТ проекти. Бизнес процесът се описва като насочен ацикличен граф
от дейности със стойности, продължителности и зависимости. Проектите кандидати
се оценяват по няколко критерия чрез два независими метода, а изборът на
портфейл се формулира като задача на целочисленото линейно програмиране с
ограничения по бюджет, по капацитет и по предхождане. Към препоръката се прилага
анализ на чувствителността, който отделя устойчивата ѝ част от частта, която се
мени заедно с допусканията.

\textbf{Предимства на решението.} Допусканията са в текстов файл, който се
версионира и сравнява, така че препоръката е възпроизводима. Оценяването и
подборът са отделени, което позволява методът за оценяване да бъде сменен, без
да се пипа подборът. Лакомата евристика, изпълнявана успоредно с решателя, дава
евтина проверка, че моделът не е сгрешен.

\textbf{Недостатъци и ограничения.} Системата приема оценките по критериите за
дадени и не проверява как са получени; качеството на препоръката не може да
надхвърли качеството на този вход. Формулировката е статична и не отчита
възможността даден проект да бъде отложен вместо отхвърлен. Обхватът спира при
одобряването на портфейла и не покрива проследяването на изпълнението.

\TODO{обобщение на действителните експериментални резултати; пише се след
Раздел~\ref{sec:results} и не го изпреварва}

\textbf{Посоки за по-нататъшна работа.} Първата е добавянето на времево
измерение, при което решението не е дали един проект да бъде изпълнен, а в кой
период да започне. Втората е вграждането на несигурността в оценките, чрез
интервали вместо точкови стойности. Третата е свързване с работеща система за
планиране на ресурсите на предприятието, така че разчетните стойности и
наличният капацитет да се четат оттам, а не да се въвеждат ръчно.
